\documentclass[12pt]{article}
\usepackage[margin=1in]{geometry}
\usepackage{amsmath,amsthm,amssymb,scrextend}
\usepackage{array}
\setlength{\parindent}{0pt}
\begin{document}
\newcommand{\ds}{\displaystyle}

\textbf{(i)}

For $L_a$,
$$y-0=\frac{(1-a)-0}{0-a}(x-a)$$
$$y=-\frac{1-a}{a}x+1-a$$

For $L_b$,
$$y=-\frac{1-b}{b}x+1-b$$


$$y=-\frac{1-a}{a}x+1-a=-\frac{1-b}{b}x+1-b$$

$$\frac{b-a}{ab}x=b-a$$

$\ds \because a\neq b,\quad \therefore x=ab$

$$y=(1-a)(1-b)$$
Points of intersection:  $ \ds (ab,\ (1-a)(1-b))$

\bigskip\hrule\bigskip
\textbf{(ii)}

$$b\to a\implies (ab,\ (1-a)(1-b))\to (a^2,\ (1-a)^2)$$

Let $\ds x=a^2\quad (0<a<1)$, 

$$y=(1-a)^2=(1-\sqrt{x})^2$$

\bigskip\hrule\bigskip
\textbf{(iii)}

$$C:\ y=(1-\sqrt{x})^2=1-2\sqrt{x}+x$$

$$\frac{dy}{dx}=1-\frac1{\sqrt{x}}=1-\frac1a=-\frac{1-a}{a}$$

\begin{align*}
y-(1-a)^2&=-\frac{1-a}{a}(x-a^2)\\
y&=-\frac{1-a}{a}x+1-a
\end{align*}

$$\therefore L_a \text{ is tangent to } C \text{ at } (a^2,(1-a)^2)$$

\bigskip
\hrule
\begin{center}
    This solution was submitted by \textbf{Jon}\\[4pt]
    Published and refined by \textbf{Jon}\\[8pt]
    \textit{Last updated: July 1, 2026}
\end{center}

\end{document}